\section{Discussion}
\label{sect:discussion}

\subsection{What the Synthetic Experiment Establishes}

The central result is not merely that a two-Gaussian profile differs from its flux-weighted centroid. Rather, the joint inference shows a broad controlled regime in which the profile is still single-peaked, yet the peak coordinate is sufficiently non-photocentric that a full eleven-parameter orbit fit cannot remove the discrepancy without degrading the joint likelihood.

The light-fraction dependence is physically interpretable within the restricted surrogate. The effect vanishes toward a dark secondary because the second light profile disappears, and it vanishes toward exact equal-light symmetry in the single-peak limit because the central maximum and photocentre coincide. Intermediate light fractions maximize the displacement of the peak from the flux-weighted position. The fact that $\betag=0.25$ is the maximum at all 230 native sampled sky/separation combinations makes this behavior robust within the tested orbit and grid rather than a feature of one scan geometry.

The full-sky experiment also demonstrates why a single representative sky position is insufficient. The same binary parameters can produce order-of-magnitude changes in native $\dchi$ because Gaia's nominal cadence and directional sampling vary over the sky. The matched-$N$ experiment then shows that most of this large-scale variation is attributable to observation count: after all positions are reduced to 53 measurements, the correlation with native $N$ vanishes and the matched/native discrepancy ratio closely tracks the retained fraction. Residual correlations with scan-angle entropy and maximum gap show that geometry remains a secondary contributor.

\subsection{What the Experiment Does Not Establish}

Several limitations are deliberate and define the next tests.

First, the Gaussian width is a research scale for the baseline experiment. Neither $\sigma=50$ mas nor any sampled $\asig$ value is a calibrated Gaia resolution scale. The quantitative transition locations cannot be translated directly into angular separation criteria for real Gaia data.

Second, the baseline model uses the global maximum of a two-equal-width-Gaussian one-dimensional profile. This is a restricted case of the blended-source problem treated more generally by \citet{Penoyre2026}, who include an elongated PSF through an orientation-dependent effective width. The published correction to that work changes signs in several equations, and the corrected expressions must be used in the next implementation. Gaia itself uses a calibrated PLSF whose DR4 treatment includes drift-scan effects and dependences on source colour and focal-plane position \citep{Rowell2026}. Window assignment, gating, background, source matching, detection, and calibration effects are absent from the baseline surrogate.

Third, the schedules are nominal GOST-derived mission trajectories. They encode sky-position-dependent times and scan directions but not the exact source-level set of successful observations. The native full-sky study should therefore be interpreted as a controlled mission-geometry experiment, not a prediction of the number of usable epochs for a specific star.

Fourth, the synthetic study holds one physical orbit, noise prescription and external-data design fixed. Equation~(\ref{eq:scaling}) should not be generalized into a universal Gaia law $\dchi\propto N$ without tests over period, eccentricity, inclination, mass ratio, magnitude/noise, external astrometric precision, and orbit/scan phase relationships.

Finally, the frozen quantitative results use the deterministic eleven-parameter baseline fitter. The repository now also contains absolute-astrometric terms, a validation-grade posterior sampler, and response-misspecification controls, but those additions were made after the frozen experiments and are not retroactively part of the 23\,000- and 11\,040-fit result sets. Real target analyses will require a production posterior engine, the Gaia spacecraft ephemeris, and calibrated response/nuisance models.

\subsection{The Literature Gap After the Deeper Audit}

The candidate contribution has to be stated at the intersection of existing methods rather than as a new ingredient in isolation. Combined orbit inference from absolute astrometry, independent relative astrometry, and component-labelled radial velocities is established by BINARYS and related tools \citep{Brandt2021,Leclerc2023}. Indeed, \citet{Chevalier2023} already combined Gaia DR3 astrometric solutions with SB2 data to derive individual stellar masses and luminosities. Component-mass inference from unresolved Gaia astrometric orbits and photometry is also established \citep{BailerJones2026}. Neither Gaia+SB2 nor Gaia-derived component masses are novelty claims available to this work.

BINARYS nevertheless provides an unusually direct motivation for the remaining measurement problem. It models the non-trivial Hipparcos transit response when secondary light matters, but \citet{Leclerc2023} explicitly limit the analogous luminous partially resolved Gaia case because the flux contamination/LSF fitting cannot be reconstructed from the Gaia catalogue products then available. The published Gaia DR3 NSS processing likewise discarded partially resolved doubles before its orbital pipeline \citep{Halbwachs2023}. Our target regime is therefore a documented limitation of existing stellar processing rather than a regime invented for the synthetic experiment.

A physical nonlinear blended-source response is itself established. \texttt{gaiamock} implements the equal-width peak response reproduced by our baseline \citep{ElBadry2024}; its public source also contains a forward predictor that combines that response with a primary-star RV curve. Therefore even response-aware stellar Gaia+RV forward modelling cannot be claimed as new. The audited public path did not reveal the full target-level inverse combination pursued here, but absence from one code path is not proof of absence in the literature.

The Arecibo analysis removes another broader priority claim. \citet{Liu2024} fitted a binary asteroid using Gaia FPR astrometry with an unresolved/partially-resolved/resolved response and external occultation relative positions. Thus a partially resolved Gaia response has already been used inside orbital inference. That analysis is sequential, is not a stellar SB2, and has no pair of stellar RV curves; those distinctions define the remaining intersection rather than making the underlying response physics new.

Bayesian Gaia epoch-astrometry plus RV inference is established by \citet{Baycroft2026}. Public Octofitter code additionally provides epoch-level Gaia DR4 likelihood infrastructure alongside relative-astrometry and RV modules. In the specific public Gaia likelihood paths audited for this project, we did not identify a physical marginal-resolution blended-profile operator. The remaining \emph{candidate} gap is consequently very narrow: a stellar target-level inference in which that operator is fitted inside the Gaia epoch likelihood while the same dynamical model is simultaneously constrained by independent resolved relative astrometry and \emph{both} SB2 velocity curves, with response assumptions propagated into $M_1$, $M_2$, parallax, and light ratio. We make no priority claim for this exact intersection.

More importantly, the stronger scientific question is independent of software priority: \emph{how much dynamical-mass and parallax bias is caused by treating marginal-resolution Gaia measurements as an ordinary photocentre, and what fidelity of response model is required to recover calibrated inference?} That question is the reason to continue the project even if another implementation eventually occupies the same capability intersection.

\subsection{Response Fidelity as the Decisive Test}

The matched-surrogate baseline gives the resolution-aware model an intentional advantage: the response used in the fit is exactly the response used to generate the Gaia-like data. The next experiment must remove that advantage. We therefore define a response hierarchy on the same synthetic observations: $M_0$ is the ordinary photocentre, $M_1$ is the equal-width Lindegren/\texttt{gaiamock}-family peak response, and $M_2$ is an orientation-dependent response based on the corrected effective-width treatment of \citet{Penoyre2026}. A later $M_3$ should allow response parameters to vary as nuisance parameters rather than asserting them.

The principal diagnostics should then be the fractional biases
\[
\frac{\widehat M_1-M_{1,\rm true}}{M_{1,\rm true}},\qquad
\frac{\widehat M_2-M_{2,\rm true}}{M_{2,\rm true}},\qquad
\frac{\widehat\varpi-\varpi_{\rm true}}{\varpi_{\rm true}},
\]
and the empirical coverage of nominal 68 and 95 per cent posterior intervals. $\dchi$ remains a useful model-mismatch diagnostic, but a large residual preference is scientifically less important than whether the inferred masses are accurate and the posterior intervals are calibrated.

The combined visual-SB2 design also creates an identifiability opportunity. The two RV amplitudes strongly constrain the mass ratio, while resolved astrometry constrains the relative orbit. Such a system can therefore provide an externally constrained projected component separation at each Gaia epoch. In principle, the remaining mapping from that separation to the measured Gaia coordinate can be used to constrain response nuisance parameters. An ensemble of well-measured visual SB2s may therefore serve as an empirical calibrator of Gaia's close-pair response rather than merely as targets whose masses require correction.

\subsection{Gaia DR4 Opportunity and Overlap Risk}

The official Gaia DR4 expected-content table already lists provisional products directly relevant to this programme, including \texttt{epoch\_astrometry}, \texttt{epoch\_image}, and \texttt{rvs\_epoch\_data\_double}. The table is explicitly under development and subject to processing and validation changes, and the detailed DR4 processing documentation/final data model are not yet released. These are therefore official expected product names, not guessed interfaces.

This is simultaneously the main opportunity and the largest remaining novelty blind spot. The public material reviewed here does not establish whether the eventual DPAC NSS likelihood internally uses a physical marginal-resolution PLSF/image response for orbital inference. Until the DR4 processing papers and documentation are available, we should claim neither that DPAC does nor that it does not implement such an inversion.

The planned \texttt{epoch\_image} product points to the natural long-term extension. Once a blended profile becomes genuinely multi-peaked, collapsing a transit to one astrometric coordinate is no longer a robust measurement model. An image/sample-domain likelihood would model the CCD samples directly and eliminate the need for an artificial single-coordinate continuation through that regime.

Until suitable epoch/image products are public, DR3 catalogue diagnostics can provide only indirect consistency tests. Quantities such as \texttt{ipd\_frac\_multi\_peak} and \texttt{ipd\_gof\_harmonic\_amplitude} are qualitatively related to partially resolved structure, but they include Gaia detection, windowing, calibration, and IPD behaviour and should not be equated directly with the mode count or residuals of the present surrogate.

\section{Conclusions}
\label{sect:conclusion}

We developed a physically constrained joint orbit prototype for resolved relative astrometry, SB2 radial velocities, and Gaia-like along-scan measurements with a resolution-aware single-peak response. The main conclusions are:

\begin{enumerate}
\item The baseline equal-width response recovers the ordinary photocentre in the unresolved limit but admits a distinct single-peak, non-photocentric regime before mode splitting. The response is not new: it reproduces the published \texttt{gaiamock} response \citep{ElBadry2024} and is a restricted case of the broader treatment of \citet{Penoyre2026}.
\item Neither combined Gaia/relative-astrometry/SB2 fitting nor Gaia-based component-mass inference is new \citep{Leclerc2023,Chevalier2023,BailerJones2026}; nor is response-aware orbital inference in general, as demonstrated for Arecibo \citep{Liu2024}.
\item The surviving candidate methodological gap is the specific stellar inverse combination: a physical marginal-resolution operator inside an epoch Gaia likelihood, fitted simultaneously with independent resolved relative astrometry and both SB2 velocity curves and propagated to individual masses. We make no priority claim for this intersection.
\item Explicit orbit conventions and single/multi-peak validity checks are essential; internally self-consistent synthetic recovery alone is insufficient to detect orientation-convention errors.
\item Across 23\,000 native full-sky baseline fits, $\betag=0.25$ gives the largest median photocentre--RAA discrepancy at all 230 tested sky/separation combinations.
\item At $\betag=0.25$ and $\asig=1$, the native sky-median $\dchi$ is 2195.39; matching every sky position to 53 observations reduces it to 1180.94 and removes the native correlation with transit count.
\item The matched/native discrepancy ratio tracks the retained-transit fraction with $\rho_{\rm S}=0.923$, while residual scan-angle dependence remains.
\item These values characterize a restricted matched-surrogate synthetic design; they are not physical Gaia resolution thresholds, empirical Gaia calibration results, or evidence that the same advantage survives realistic response misspecification.
\item The decisive next question is therefore response fidelity: whether progressively more realistic response models improve mass/parallax bias and posterior coverage when the injected measurement physics is deliberately not identical to the fitted model.
\end{enumerate}

The longer-term observational programme is to use DR4 epoch astrometry and, eventually, epoch images to test the same hierarchy on real visual SB2s and to determine whether strong external orbital constraints can empirically calibrate the close-pair astrometric response.

\section*{Data and Code Availability}

The implementation, validation tests, analysis scripts, compact frozen result tables, and manuscript figures are versioned in the accompanying \texttt{RAA\_Orbit\_Model} repository. The large raw synthetic CSV files are generated by the documented scripts; the committed frozen tables contain the numerical summaries used in this manuscript. Exact nominal scan schedules are archived by the experiment runner so that the same directional sampling can be re-used without re-querying the scanning-law provider.

\begin{acknowledgements}
The author acknowledges the open Gaia mission documentation and the public software and literature resources used to construct the controlled scanning-law experiments.
\end{acknowledgements}
